\chapter{Model Architecture}

\section{Face Detection Models}
Face detection is the initial stage in our attendance pipeline. To ensure the system operates reliably across different hardware environments, we implement a multi-model fallback strategy. The detection layer supports three face detection architectures:

\begin{enumerate}
    \item \textbf{YOLOv8-face:} A state-of-the-art anchor-free object detector specialized for face detection and landmark localization. Employs a backbone featuring C2f (CSP Bottleneck with 2 convolutions) blocks for multi-scale feature extraction. YOLOv8-face is highly robust against facial occlusions, varied head poses, and diverse lighting conditions. It is the primary face detector in our system.
    \item \textbf{MTCNN (Multi-task Cascaded Convolutional Networks):} A three-stage cascaded deep learning framework containing:
    \begin{itemize}
        \item \textit{Proposal Network (P-Net):} A fully convolutional network that generates candidate bounding boxes.
        \item \textit{Refinement Network (R-Net):} A CNN that rejects false candidates and calibrates bounding boxes.
        \item \textit{Output Network (O-Net):} A more complex CNN that performs final refinement and extracts five facial landmarks (eyes, nose, and mouth corners).
    \end{itemize}
    MTCNN serves as the primary fallback when YOLOv8-face dependencies are unavailable.
    \item \textbf{Haar Cascade:} OpenCV's Viola-Jones face detector \cite{viola2001rapid}. It utilizes Haar-like features and a cascade of AdaBoost classifiers. It is highly lightweight and runs on extremely low-power CPUs without GPU or deep learning libraries (e.g., PyTorch), serving as the final fallback for development environments.
\end{enumerate}

This defensive fallback configuration allows the pipeline to degrade gracefully if model files or PyTorch libraries are missing, guaranteeing continuous uptime in production.

\subsection{YOLOv8-face Architecture Details}
YOLOv8-face represents the state-of-the-art in face detection. Unlike classical anchor-based detectors, it is an anchor-free model, directly predicting the center offset and dimensions of bounding boxes. Its backbone is based on a modified CSPDarknet53, which incorporates the C2f module to merge high-resolution spatial details with low-resolution semantic features. 

The detection head is split into two branches:
\begin{itemize}
    \item \textbf{Bounding Box Regression Branch:} Predicts the box coordinates $[x, y, w, h]$ and regress the five key facial landmarks: left eye, right eye, nose tip, left mouth corner, and right mouth corner. Bounding box coordinates are optimized using Complete Intersection over Union (CIoU) loss.
    \item \textbf{Classification Branch:} Predicts the probability of a face presence in the candidate region using Binary Cross-Entropy (BCE) loss.
\end{itemize}
The landmark predictions allow the system to perform affine transformations on the face crops, aligning the eyes horizontally before embedding extraction, which increases face recognition accuracy.

\subsection{MTCNN Cascade Details}
MTCNN operates in three sequential convolutional stages to optimize speed and accuracy:
\begin{enumerate}
    \item \textbf{P-Net ($12\times12$ input):} A light, fully-convolutional network that scans the input image pyramid to generate candidate windows. Boundary box regression vectors are used to calibrate candidates, followed by Non-Maximum Suppression (NMS) with a threshold of 0.70 to merge overlapping regions.
    \item \textbf{R-Net ($24\times24$ input):} Candidate patches from P-Net are fed into R-Net, which uses a fully-connected layer to reject false alarms. A second NMS step (threshold 0.70) is applied.
    \item \textbf{O-Net ($48\times48$ input):} The final stage identifies face patches in detail, outputting the bounding box coordinates, a face probability score, and the coordinates of five landmarks. Overlapping regions are merged using a strict NMS threshold of 0.60.
\end{enumerate}

\subsection{Haar Cascade (Viola-Jones) Concept}
Haar Cascade uses Haar-like features (rectangular features that compute intensity differences between adjacent regions). To compute these features in constant time $O(1)$, an \textbf{Integral Image} (or summed-area table) is created. 
The integral image $II(x, y)$ contains the sum of pixel intensities above and to the left of $(x, y)$:
\begin{equation}
    II(x, y) = \sum_{x' \le x, y' \le y} I(x', y')
\end{equation}
AdaBoost is used to select the most discriminative features and construct a cascade of classifiers. When checking a region, early cascade stages quickly reject background patches, allowing the detector to run fast on low-power CPUs. However, it is sensitive to head rotation, shadows, and variations in skin tone.

\section{Face Recognition Model}
Once a face bounding box is detected, it is cropped and passed to the face recognition module to extract biometric features.

\subsection{InceptionResnetV1 (FaceNet)}
We utilize the \textit{InceptionResnetV1} architecture, trained using a Triplet Loss function to output 512-dimensional face embeddings \cite{schroff2015facenet}. The network combines the multi-scale feature representation of Inception blocks with the optimization stability of residual connections (skip connections). The residual path prevents gradient vanishing during deep backpropagation, allowing the model to learn fine-grained facial identity features.

\subsection{Embedding Verification and Cosine Similarity}
Let $A \in \mathbb{R}^{512}$ be the reference embedding stored in the database, and $B \in \mathbb{R}^{512}$ be the query embedding extracted from the live camera frame. The similarity between the two face representations is calculated using the Cosine Similarity:
\begin{equation}
    \text{Similarity}(A, B) = \frac{A \cdot B}{\|A\|_2 \|B\|_2}
\end{equation}
Since all embeddings are $L_2$ normalized to unit length during preprocessing ($\|A\|_2 = \|B\|_2 = 1.0$), the cosine similarity simplifies directly to the dot product:
\begin{equation}
    \text{Similarity}(A, B) = A \cdot B = \sum_{i=1}^{512} a_i b_i
\end{equation}
A match is declared if the similarity score exceeds a predefined threshold $\theta_{\text{verify}}$, which is calibrated to minimize false registrations.

\subsection{Triplet Loss Optimization}
The InceptionResnetV1 model is optimized using Triplet Loss, which maps face images directly into a compact Euclidean space where distances correspond to facial similarity. The training is conducted using triplets:
\begin{itemize}
    \item \textbf{Anchor Image ($x^a$):} Image of a specific person.
    \item \textbf{Positive Image ($x^p$):} Another image of the same person.
    \item \textbf{Negative Image ($x^n$):} Image of a different person.
\end{itemize}
The triplet loss function minimizes the distance between the anchor and the positive image, while maximizing the distance between the anchor and the negative image by a margin $\alpha$:
\begin{equation}
    \mathcal{L} = \sum_{i=1}^N \max\left(0, \|f(x_i^a) - f(x_i^p)\|_2^2 - \|f(x_i^a) - f(x_i^n)\|_2^2 + \alpha\right)
\end{equation}
where $f(x)$ represents the 512-dimensional embedding projected by the network.

\section{Facial Expression Recognition Model}
To implement emotion-gated check-in (requiring employees to smile to log attendance), we deploy a lightweight Convolutional Neural Network (CNN) model based on the \textit{Mini-Xception} architecture.

\subsection{Mini-Xception Architecture}
Mini-Xception is a compact neural network inspired by the Google Xception design. Standard convolutional layers are replaced with \textbf{Depthwise Separable Convolutions} to minimize parameter count and floating-point operations.

A depthwise separable convolution divides a standard convolution into two separate steps:
\begin{itemize}
    \item \textbf{Depthwise Convolution:} Applies a single spatial filter per input channel.
    \item \textbf{Pointwise Convolution:} Applies a $1 \times 1$ convolution to project the output channels into a new dimension.
\end{itemize}
The computational complexity comparison is:
\begin{itemize}
    \item Standard Convolution:
    \begin{equation}
        C_{\text{std}} = H_k \cdot W_k \cdot C_{\text{in}} \cdot C_{\text{out}} \cdot H_{\text{out}} \cdot W_{\text{out}}
    \end{equation}
    \item Depthwise Separable Convolution:
    \begin{equation}
        C_{\text{sep}} = (H_k \cdot W_k \cdot C_{\text{in}} \cdot H_{\text{out}} \cdot W_{\text{out}}) + (1 \cdot 1 \cdot C_{\text{in}} \cdot C_{\text{out}} \cdot H_{\text{out}} \cdot W_{\text{out}})
    \end{equation}
\end{itemize}
The reduction ratio is:
\begin{equation}
    \text{Ratio} = \frac{C_{\text{sep}}}{C_{\text{std}}} = \frac{1}{C_{\text{out}}} + \frac{1}{H_k \cdot W_k}
\end{equation}
For standard $3 \times 3$ kernels ($H_k = W_k = 3$), this ratio is approximately:
\begin{equation}
    \text{Ratio} \approx \frac{1}{9} + \frac{1}{C_{\text{out}}}
\end{equation}
which represents a saving of ~89\% in floating-point operations.

Additionally, the architecture utilizes residual connections, Batch Normalization, and rectified linear unit (ReLU) activations. Rather than using large, fully-connected (dense) layers at the output, the network uses \textbf{Global Average Pooling (GAP)} immediately before the final Softmax layer. This reduces parameters significantly, resulting in a model file size under 5MB.

\begin{table}[H]
\centering
\caption{Mini-Xception network layers configuration.}
\label{tab:mini-xception-layers}
\begin{tabular}{llcc}
\toprule
\textbf{Layer Type} & \textbf{Specification / Kernel} & \textbf{Input Shape} & \textbf{Output Channels} \\
\midrule
Input & Gray Scale Image & $1 \times 64 \times 64$ & 1 \\
Conv2d & $3 \times 3$, padding=1 & $1 \times 64 \times 64$ & 8 \\
BatchNorm + ReLU & -- & $8 \times 64 \times 64$ & 8 \\
Conv2d & $3 \times 3$, padding=1 & $8 \times 64 \times 64$ & 8 \\
BatchNorm + ReLU & -- & $8 \times 64 \times 64$ & 8 \\
\midrule
\textbf{Residual Block 1} & 2 $\times$ Depthwise SepConv & $8 \times 64 \times 64$ & 16 \\
Max Pooling & $2 \times 2$, stride=2 & $16 \times 64 \times 64$ & 16 \\
\midrule
\textbf{Residual Block 2} & 2 $\times$ Depthwise SepConv & $16 \times 32 \times 32$ & 32 \\
Max Pooling & $2 \times 2$, stride=2 & $32 \times 32 \times 32$ & 32 \\
\midrule
\textbf{Residual Block 3} & 2 $\times$ Depthwise SepConv & $32 \times 16 \times 16$ & 64 \\
Max Pooling & $2 \times 2$, stride=2 & $64 \times 16 \times 16$ & 64 \\
\midrule
\textbf{Residual Block 4} & 2 $\times$ Depthwise SepConv & $64 \times 8 \times 8$ & 128 \\
Max Pooling & $2 \times 2$, stride=2 & $128 \times 8 \times 8$ & 128 \\
\midrule
Global Avg Pool & Average spatial dimension & $128 \times 4 \times 4$ & 128 \\
Flatten / Linear & Dense projection & 128 & 8 \\
\bottomrule
\end{tabular}
\end{table}
This compact architecture is exported as an ONNX model, allowing OpenCV DNN to load and execute inference on CPUs in under 3 milliseconds per face, which is critical for real-time video feeds.
